% !TEX root = ../main.tex
% ---------------------------------------------------------------------
% The Predictive Geometry of Reasoning — Movement 1 (Structure) section.
% DRAFT. Side project of the thesis; canonical content source is
% ../PREDICTIVE_GEOMETRY.md. Voice + rules per thesis/_planning/STYLE.md.
%
% NUMBER EMBARGO (STYLE.md §5): no measured numerals appear here. Pending
% results use prose ("pending the pilot gate") or the thesis embargo macros.
% The prompt asked for a \TODO{} macro for unmeasured quantities; the thesis
% preamble defines \gateZero{...} (results owed to a gate) and \todoT{...}
% (decisions owed to Tony) instead, so those are used here. If this file is
% compiled standalone rather than via thesis/main.tex, define them, e.g.:
%   \providecommand{\gateZero}[1]{\textbf{[GATE]}\emph{#1}}
%   \providecommand{\todoT}[1]{\textbf{[TODO]}#1}
%
% CITATION KEYS (see predictive_geometry_keys.txt for every key + arXiv id):
%   Reused, already in thesis/references.bib:
%     sun2026trajectories (2604.05655), zhou2025flowing (2510.09782),
%     raval2026curveball (curved steering), venhoff2025reasoning (steering).
%   NEW, NOT yet in the bib — must be added before compile, or downgraded to
%   \needcite{...} per STYLE.md §6: ssp2026smoothness (2604.18464),
%   hao2024coconut (2412.06769), llmjepa2025 (2509.14252),
%   phasespace2024 (2410.04415), assael2025jepa (I-JEPA), bardes2025vjepa2,
%   balestriero2025lejepa, zbontar2021barlow, asm2026 (ASM, ICLR 2026).
% ---------------------------------------------------------------------
\section{The predictive geometry of reasoning}\label{sec:predictive-geometry}

The previous chapter closed one rung of the structural ladder and left a door
deliberately open. Per-behaviour curvature, measured on the cloud of mean-pooled
span activations, was tested and not found: the bend the early runs reported
dissolved under the chain control, exposed as the local clustering of
neighbouring sentences from the same chain rather than a geometry of the
behaviour (\Cref{ch:geometry}). The verdict came with a caveat that pointed
forward as much as it settled anything. Mean-pooling each span to a single point
discards the very thing a curvature claim might most naturally concern, the path
a chain traces as a behaviour unfolds, and so \Cref{ch:geometry} relocated
curvature, as an open question, to the trajectory: if it lives anywhere in this
data it lives there, left to future work. The conceptual chapter had already
argued for the same locus from a different direction. Reading the AlphaProof
Nexus basic-agent finding at its finest grain, \Cref{ch:reasoning} committed to
locus-in-process, the position that it is the unfolding trajectory and not any
static activation point that bears the reasoning, and observed that the natural
future experiments are therefore trajectory-based. This section builds the
experiment those two commitments name. It is the trajectory-level realisation of
a question the rest of Movement~1 could pose but not, with a static apparatus,
answer.

The move is to stop describing geometry and start predicting it. Everything in
the structural programme so far has read shape off a fixed set of points: the
intrinsic dimension of a behaviour's cloud, the curvature of that cloud, the
subspace it spans. A trajectory invites a different question. Given the
embeddings of the reasoning steps a chain has produced so far, $x_1$ through
$x_t$, can we predict the next one, $x_{t+1}$? The object of interest is then not
the static cloud but the residual $r_t = x_{t+1} - f(x_{\le t})$, the part of the
next step a learned forward model could not anticipate. The wager is that this
residual is informative where the static geometry is silent: that the predictor
breaks, with sharp residuals, exactly at the branch points and backtracks where
reasoning changes direction, and that the structure of those breaks separates
chains that reach a correct answer from chains that do not. Descriptive geometry
asks what reasoning representations look like at rest. The predictive reframing
asks how they move, and treats the predictability of that motion as the signal.

The apparatus this calls for already exists, under the name of the
joint-embedding predictive architecture. The JEPA programme \citep{assael2025jepa,
bardes2025vjepa2} is built to predict future states in a learned latent space
while discarding the unpredictable detail a generative model would waste capacity
reconstructing; LeJEPA \citep{balestriero2025lejepa} adds an isotropy regulariser
that prevents the learned space from collapsing to a point. That is precisely the
machine the trajectory question wants. We do not need to reconstruct the next
reasoning sentence token by token, only to predict where in latent space it will
land and to read the error when it lands elsewhere. Naming JEPA as the apparatus
fixes the register of the whole section: the thesis is not proposing a new
architecture but borrowing a mature one to answer a question its own static
geometry left open.

Four hypotheses organise the work, and they are graded so that the cheaper claims
gate the dearer ones. The keystone is functional. \textbf{H1} holds that the
predictability of a reasoning trajectory --- the magnitude of its residuals,
their growth across the chain, the churn in their direction --- predicts whether
the chain is correct, and predicts it better than token-level log-likelihood, raw
curvature, or chain length and difficulty alone. \textbf{H2} is representational:
that a \emph{learned} predictor recovers correctness and behaviour structure that
neither the raw hidden states nor a random projection of them expose, so that the
signal is in the learning and not merely in the embeddings. \textbf{H3} is
geometric and ties the residual back to the static picture: that residual spikes
localise at points the static geometry would call special --- high local
curvature, transitions between subspaces --- and that these are where correct and
incorrect trajectories part company once difficulty is controlled, the
trajectory-level home \Cref{ch:geometry} reserved for a curvature it would not
assert of the behaviour cloud. \textbf{H4} is the causal apex, and it is offered
as confirmation rather than novelty: that steering the model along the
\emph{predicted} next-step direction shifts its reasoning more than steering along
a random or mean direction, closing the loop from a passive diagnostic to an
intervention. The four are nested. If the residual carries no correctness signal
at all, the geometric and causal questions do not arise; H1 is the gate the rest
must pass.

The method is an escalation ladder, and each rung is run only if it beats the rung
below it and survives two null tests. The discipline is inherited wholesale from
the rest of the thesis. Every split is chain-grouped, so a chain never trains and
tests at once and the effective sample size is the chain count and not the far
larger step count (the chain confound the apparatus chapter names CF-2); every
positive claim must clear both a chain-stratified label permutation, which asks
whether the signal exceeds chance given the composition of chains, and a
within-chain step-order shuffle, which asks whether it needs the genuine temporal
order or is a static-cloud artefact wearing a trajectory's clothes. Rung~0 is
predictor-free. It runs the existing Frenet-curvature diagnostic on the raw
trajectories of difficulty-matched correct and incorrect chains, to learn whether
\emph{any} trajectory geometry separates correctness before we build a predictor
at all. Rung~1 is linear and, by design, carries no anti-collapse regulariser: a
ridge model predicts the displacement $x_{t+1}-x_t$ rather than the absolute next
state, so the residual measures only what the predictor improves on persistence,
and because nothing regularises the latent, reading geometry off the residual is
not circular. Rung~2 introduces a small JEPA --- a compact causal predictor over
step embeddings, with an anti-collapse term of the LeJEPA or Barlow-Twins kind
\citep{balestriero2025lejepa, zbontar2021barlow} --- and it is pursued only to
\emph{beat} Rung~1, and only ever for relative, functional claims, never for an
absolute intrinsic dimension or isotropy of a latent its own regulariser has
shaped. Rung~3 is the apex: goal-conditioned latent rollout with planning, and
causal steering along the predicted direction through the steering apparatus of
\Cref{ch:apparatus}. The ladder is built so that a negative at any rung is
informative and a positive is hard-won.

Placing this work against its nearest neighbours matters more here than usual,
because the surrounding literature is dense and the genuinely novel atom is
narrow. The descriptive geometry of reasoning trajectories has been mapped from
several directions: a body of work treats reasoning as flowing logics in
representation space \citep{zhou2025flowing}, and a phase-space reading studies the
dynamics of the hidden state directly \citep{phasespace2024}. Predicting the next
reasoning step in a latent space is also not new in itself --- COCONUT
\citep{hao2024coconut} reasons in a continuous latent rather than in tokens, and
LLM-JEPA \citep{llmjepa2025} trains a paired-view predictor over text --- but both
are generative or representational ends in themselves, not error signals read for
diagnosis. Steering along a predicted or idealised next-state direction has been
proposed as well \citep{asm2026}, which is exactly why H4 is framed as
confirmation. Two results sit closest, and the wedge runs between them. A static
activation-difference probe reaches high accuracy at telling correct from
incorrect reasoning, but from fixed representations rather than from any model of
how they evolve \citep{sun2026trajectories}; and a study of trajectory smoothness
reports that smoothness does \emph{not} encode correctness, returning an
area-under-curve near chance, and explicitly declines to build a
latent-prediction feedback loop \citep{ssp2026smoothness}. Our quantity is neither
of theirs. It is the residual geometry of a \emph{learned forward predictor} ---
the error of a model trained to anticipate the next step, read as a correctness
and branch-point signal. No prior work trains such a predictor over reasoning
steps and treats its failure as the diagnostic, and the quantity aims squarely at
the open question the smoothness study left negative. That is the whole of the
novelty claim, and it is deliberately small: a third statistic in a space whose
two obvious endpoints are already occupied.

What the pilot will report is a single gate, stated now without numbers because no
number is yet earned. The pilot runs Rung~0 and Rung~1 on the standing R1-1.5B
corpus and asks one question: does the residual geometry of the linear predictor
predict chain-level correctness, scored as a chain-grouped area under the curve,
better than the baselines that ought to be hard to beat --- token-level
log-likelihood, raw curvature, and a length-and-difficulty model --- and does it
do so while \emph{both} nulls reject? An early, pre-label read already tempers
expectation: a cross-chain linear next-step map appears to underperform simple
persistence, which would make Rung~1 a null and hand the work to Rung~2, sharpening
rather than weakening the argument against the smoothness result. The honest
spread of outcomes is therefore wide. A clean positive, with the residual beating
every baseline on difficulty-matched chains and both nulls rejecting, would move
the thesis's claim that reasoning is a geometric process from a descriptive
statement about static representations to a dynamical one about how reasoning
unfolds. An informative null --- Rung~1 at chance but a nonlinear Rung~2 clearing
the bar --- would say the signal is real but nonlinear, that capacity is what the
trajectory question needs. And a hard null, neither rung beating the baselines
with the nulls intact, would be a clean negative that confirms and extends the
smoothness result: trajectory smoothness, static probes read dynamically, and now
the residual geometry of a learned predictor would all have failed to find
correctness in the shape of the path, a genuine and reportable limit on how much
of reasoning its geometry carries. The reframe is worth running for the same
reason \Cref{ch:geometry} was worth waiting for: each of the three outcomes
teaches, and which one obtains is not yet known.

\gateZero{Pilot gate: chain-grouped correctness AUC for the Rung~0 curvature
diagnostic and the Rung~1 linear residual, against the token-NLL, raw-curvature,
and length-plus-difficulty baselines, each reported with its chain-stratified
label-permutation null and its within-chain step-order-shuffle null. No value is
quoted until the pilot has run on difficulty-matched chains.}

\todoT{Decide whether the predictive-geometry section ships inside
\Cref{ch:geometry} as its closing trajectory move or as a standalone Movement~1
section; resolve the nine new citation keys (predictive\_geometry\_keys.txt)
into thesis/references.bib, or downgrade any that cannot be verified to
\needcite. Recommended once the pilot gate has a verdict to anchor the framing.}
